% ============================================================================
%  Speaker notes for review.tex --- one entry per slide.
%  Build:  cd slides && latexmk -pdf speaker_notes.tex
%  ⛔ No number is typed here either; macros come from paper/numbers.tex.
% ============================================================================
\documentclass[11pt,a4paper]{article}
\usepackage[margin=2.2cm]{geometry}
\usepackage[T1]{fontenc}
\usepackage{lmodern,xcolor,amsmath,booktabs,enumitem,titlesec}
\usepackage[colorlinks=true,linkcolor=blue!50!black,urlcolor=blue!50!black]{hyperref}

\definecolor{acc}{HTML}{1F3F8B}
\definecolor{warn}{HTML}{B02418}
\newcommand{\best}[1]{\textcolor{acc}{\textbf{#1}}}
\newcommand{\bad}[1]{\textcolor{warn}{#1}}
\newcommand{\ok}[1]{\textcolor{acc}{#1}}
\newcommand{\bg}{\texttt{background}}
\newcommand{\suppref}[1]{}
% =====================================================================
%  Every load-bearing number in the paper, defined ONCE.
%
%  WHY THIS FILE EXISTS. This project has caught five silent errors, and
%  the cheapest class of error left is a number that is right in the
%  results table and stale in the abstract. Macros make that impossible:
%  change it here and it changes everywhere, or it changes nowhere.
%
%  Each macro carries the file and section it came from. If a macro has
%  no citation it is not allowed in the paper.
% =====================================================================

% ---- baselines ------------------------------------------------------
\newcommand{\loveBase}{47.38}        % WEEK1 §5   reproduced, full val
\newcommand{\lovePub}{47.4}          % SegEarth-OV3, Table 1
\newcommand{\loveInstr}{47.37}       % WEEK1 §5   instrumented recompute (the gate)
\newcommand{\oemBase}{44.19}         % WEEK3 §7   384 of 500 val tiles
\newcommand{\oemPub}{42.9}           % SegEarth-OV3, published
\newcommand{\loveTiles}{1669}
\newcommand{\oemTiles}{384}

% ---- the residual ---------------------------------------------------
\newcommand{\discardPct}{29.68\%}    % WEEK1 §7.1
\newcommand{\discardPx}{323{,}184{,}908}
\newcommand{\discardOEM}{3.78\%}     % WEEK3 §7
\newcommand{\discardRatio}{3:1}      % WEEK1 §7.6  discard vs confusion
\newcommand{\mechA}{94.0\%}          % WEEK1 §7.7  below tau
\newcommand{\mechB}{6.0\%}           % WEEK1 §7.7  background won the argmax
\newcommand{\mechBpx}{19{,}378{,}177}
\newcommand{\mechBtiles}{24}
\newcommand{\mechBpots}{1.25\%}    % POTSDAM, reachability.py 4 Sep: argmax loss
\newcommand{\cnOemFloor}{0.1042}    % CONINFER_RESULTS: their OEM cache conf floor
\newcommand{\cnOemTau}{0.1}         % their published prob_thd on OpenEarthMap


% ---- the mechanism --------------------------------------------------
\newcommand{\bgShareLove}{36.1\%}    % WEEK3 §7
\newcommand{\bgShareOEM}{0.84\%}
\newcommand{\aucLove}{0.622}         % best of nine signals
\newcommand{\aucOEM}{0.913}          % conf2
\newcommand{\aucFloor}{0.53}         % empirical, size-stratified
\newcommand{\presCorrLove}{-0.750}
\newcommand{\presCorrOEM}{+0.094}
\newcommand{\catTilesLove}{198}
\newcommand{\catTilesOEM}{0}

% ---- refuted interventions -----------------------------------------
\newcommand{\tauCost}{-5.54}         % WEEK1 §7.2  tau 0.5 -> 0.1
\newcommand{\tauRatio}{1.73}         % WEEK1 §8.2  wrong per right
\newcommand{\presCost}{-11.97}       % WEEK1 §9.2b
\newcommand{\recoverLove}{+0.04}     % WEEK3 §8
\newcommand{\recoverOEM}{+2.28}      % WEEK3 §8   -- NEVER without \oemBgGain
\newcommand{\oemBgGain}{+22.67}      % WEEK3 §8.1
\newcommand{\oemRealLoss}{-2.11}     % WEEK3 §8.1
\newcommand{\oemBuildLoss}{-3.75}    % WEEK3 §8.1
\newcommand{\priorGain}{+0.2}        % WEEK3 §6   over a neighbour vote
\newcommand{\priorOracle}{+0.3}      % WEEK3 §6   with a PERFECT matrix

% ---- atomisation ----------------------------------------------------
\newcommand{\ceilCC}{72.8\%}         % WEEK3 §4
\newcommand{\ceilSLIC}{92.8\%}
\newcommand{\ceilSLICoem}{82.9\%}

% ---- the method -----------------------------------------------------
\newcommand{\methodGain}{+1.18}      % WEEK3 §9b  5-fold mean
\newcommand{\methodSD}{0.45}
\newcommand{\methodWorst}{+0.80}
\newcommand{\methodLand}{+8.30}      % 5-fold, six real classes
\newcommand{\methodBg}{-0.01}
\newcommand{\methodWater}{+6.78}
\newcommand{\methodWaterTau}{0.195}
\newcommand{\oracleBound}{+1.46}     % WEEK3 §9a  per-class oracle
\newcommand{\oracleShare}{81\%}
\newcommand{\globalOracle}{+0.04}    % WEEK3 §9a  best GLOBAL tau
\newcommand{\calibTiles}{200}        % WEEK3 §9b  every draw positive
\newcommand{\trainToVal}{-0.12}      % WEEK3 §9b  scope limit

% ---- end-to-end verification ---------------------------------------
\newcommand{\eeBase}{47.16}          % WEEK3 §9c  held-out tiles
\newcommand{\eeFit}{48.35}
\newcommand{\eeTiles}{1469}
\newcommand{\eeSpread}{0.175--0.675}

% ---- the label-free bound ------------------------------------------
\newcommand{\proxyBest}{+0.42}       % WEEK3 §9d  presence, best of eight
\newcommand{\proxyCtrl}{+0.58}       % random-control p95
\newcommand{\proxyConfRho}{+0.943}   % mean_conf vs precision
\newcommand{\proxyConfP}{0.017}
\newcommand{\proxyHeadRho}{+0.886}   % sem_inst_agree vs precision
\newcommand{\proxyHeadP}{0.033}
\newcommand{\proxyTauRho}{-0.429}    % precision vs ORACLE TAU -- the broken link
\newcommand{\proxyTauP}{0.419}
\newcommand{\nAttempts}{eleven}      % 3 threshold rules + 8 proxies

% ---- the co-occurrence prior ---------------------------------------
\newcommand{\pmiSignal}{0.574}       % ANALYSIS §4, corrected marginal
\newcommand{\pmiFloor}{0.003}
\newcommand{\pmiRatio}{216}
\newcommand{\rhoMined}{+0.757}       % WEEK3 §5  gate 1
\newcommand{\rhoCirc}{-0.257}        % circularity, retired

% ---- the confound stratification (WEEK3 §7a) -----------------------
\newcommand{\SHAREr}{42.9\%}       % LoveDA rural, catch-all share of GT
\newcommand{\SHAREu}{26.0\%}       % LoveDA urban
\newcommand{\CONFr}{25.5\%}        % rural confusability, measured
\newcommand{\CONFu}{43.3\%}        % urban -- the LOWER-share stratum is more confusable
\newcommand{\AUCr}{0.524}          % rural detection AUC (conf)
\newcommand{\AUCu}{0.730}          % urban

% ---- WEEK3_RESULTS.md 9e: per-class tau across LoveDA's domains -------------
\newcommand{\DOMr}{+2.77}          % rural within-domain gain, 5-fold
\newcommand{\DOMrsd}{0.92}
\newcommand{\DOMu}{+0.10}          % urban -- NOT distinguishable from zero (2/5 folds)
\newcommand{\DOMusd}{0.39}
\newcommand{\DOMrreal}{+18.61}     % rural, real classes aggregate
\newcommand{\DOMureal}{+4.18}      % urban, real classes aggregate -- land cover DOES gain
\newcommand{\DOMubg}{-3.51}        % urban catch-all pays for it: (4.18-3.51)/7 = +0.10
\newcommand{\DOMtaudiff}{0.227}    % mean |tau difference| between the domains
\newcommand{\DOMtaumax}{0.500}     % max, on `road' (0.725 rural vs 0.225 urban)
\newcommand{\DOMxr}{-0.40}         % mismatched -> rural, BELOW the published tau
\newcommand{\DOMxu}{-1.11}         % mismatched -> urban, BELOW the published tau
\newcommand{\DOMmr}{+2.32}         % matched -> rural, same 200-tile budget
\newcommand{\DOMmu}{+0.02}         % matched -> urban
\newcommand{\DOMpr}{+0.77}         % pooled -> rural: keeps a third of matched
\newcommand{\DOMpu}{-0.32}         % pooled -> urban

% ---- WEEK3_RESULTS.md 7b: the vocabulary intervention ----------------------
% det = max(AUC, 1-AUC). AUC is symmetric, so a raw AUC of 0.208 is a detector
% of strength 0.792 inverted, NOT an absent signal. Never quote raw AUC here.
\newcommand{\ivA}{0.794}           % A0, published vocabulary, det(conf)
\newcommand{\ivAtwo}{0.913}        % A0, det(conf2)
\newcommand{\ivShareA}{0.84\%}     % catch-all share, published
\newcommand{\ivShareD}{58.20\%}    % catch-all share, largest dose
\newcommand{\ivB}{0.582}           % faithful dose arm, det(conf)
\newcommand{\ivBtwo}{0.590}        % faithful dose arm, det(conf2)
\newcommand{\ivD}{0.710}           % arity-matched control, det(conf)
\newcommand{\ivDtwo}{0.855}        % arity-matched control, det(conf2)
\newcommand{\ivRatio}{2.5$\times$} % dose effect / control effect, conf
\newcommand{\ivRatioTwo}{5.6$\times$}  % ... conf2
\newcommand{\ivRatioWorst}{3.1$\times$} % conf2 vs the WORST control, not the matched one
\newcommand{\ivArms}{13}           % arms, all reading one cached score stack

% ---- WEEK3_RESULTS.md 7c: the reverse arm, and what bounds the mechanism ----
\newcommand{\ivRev}{0.540}         % LoveDA, catch-all removed from the VOCABULARY
\newcommand{\ivRevBase}{0.592}     % LoveDA A0 -- vocabulary change alone does nothing
\newcommand{\ivVocabMax}{0.052}    % largest move by ANY arm that leaves GT share alone
\newcommand{\ivSatShare}{35\%}     % above this, four points sit in one band
\newcommand{\ivSatLo}{0.58}
\newcommand{\ivSatHi}{0.64}

% ---- WEEK3_RESULTS.md 9f: no label-free rule for WHEN calibration pays ------
\newcommand{\crProxyRho}{+0.885}   % label-free catch-all fraction vs labelled discard, LoveDA
\newcommand{\crProxyRhoOEM}{+0.924}
\newcommand{\crLow}{$+2.13 \pm 0.18$}   % LEAST-residual stratum -- the most reliable gain
\newcommand{\crHigh}{$+3.22 \pm 3.15$}  % most-residual stratum -- worst fold -1.22
\newcommand{\crRho}{+0.400}        % LoveDA, over 4 strata
\newcommand{\crRhoOEM}{-0.500}     % OEM -- opposite sign, nothing replicates

% ---- WEEK3_RESULTS.md 9g: what the domain gap is made of --------------------
\newcommand{\cpWater}{48\%}        % water's share of the rural/urban gap
\newcommand{\cpForest}{37\%}       % forest's share
\newcommand{\cpTop}{85\%}          % the two together
\newcommand{\cpGapRho}{+0.713}     % P-R gap vs per-class delta IoU, 12 cells
\newcommand{\cpGapP}{0.013}        % permutation p
\newcommand{\cpPrecRho}{+0.168}    % precision ALONE -- explains nothing
\newcommand{\cpPrecP}{0.60}

% ---- WEEK3_RESULTS.md 9h: both metrics ------------------------------------
% ⚠️ ALWAYS the per-class MEAN (= catch-all-excluded mIoU), never the aggregate
% sum. Sections 9b/9e quote sums; +8.30 aggregate IS +1.38 here. Mixing them puts
% two numbers 6x apart in units next to each other.
\newcommand{\mrLoveFull}{+1.16}    % LoveDA, full mIoU, 5-fold held out
\newcommand{\mrLoveReal}{+1.36}    % LoveDA, catch-all-excluded mIoU
\newcommand{\mrLoveBg}{-0.02}
\newcommand{\mrOemFull}{+0.30}     % OEM, full mIoU -- looks flat
\newcommand{\mrOemReal}{+1.75}     % OEM, catch-all-excluded -- LARGER than LoveDA
\newcommand{\mrOemBg}{-11.30}
\newcommand{\mrOemLand}{+1.56}     % the real classes' contribution to full mIoU
\newcommand{\mrOemBgC}{-1.26}      % background's contribution: cancels it
\newcommand{\mrOemDistort}{3.38}   % OEM's baseline headline depression
\newcommand{\mrLoveLever}{14.3\%}  % 1/N -- the catch-all's share of the metric
\newcommand{\mrOemLever}{11.1\%}
\newcommand{\mrOemBgIoU}{17.13}
\newcommand{\methodLandMean}{+1.38}  % 9b's +8.30 aggregate, in mIoU units

% ---- CONINFER_RESULTS.md 7.1a: per-class tau on a CLIP backbone -------------
% ⚠️ Deltas measured on OUR reproduction (36.99), not their published 39.33. A
% delta is robust to a constant offset, but the text must say so.
\newcommand{\cnPub}{39.33}         % ConInfer's published LoveDA mIoU
\newcommand{\cnRepro}{36.99}       % our reproduction -- gap -2.34, full split
\newcommand{\cnGap}{2.34}          % CONINFER_RESULTS: repro below published
\newcommand{\cnTau}{0.8}           % THEIR published LoveDA threshold
\newcommand{\cnBase}{37.00}        % held-out baseline at their tau
\newcommand{\cnFit}{39.52}         % held out, per-class tau
\newcommand{\cnFull}{+2.52}
\newcommand{\cnReal}{+1.94}        % catch-all-excluded -- above SAM 3's +1.36
\newcommand{\cnBg}{+6.01}
\newcommand{\cnCV}{+2.51}          % five-fold mean
\newcommand{\cnCVsd}{0.34}
\newcommand{\cnWorst}{+2.22}       % every fold positive
\newcommand{\cnCalib}{25}          % tiles for a reliable gain -- 8x cheaper here
% ConInfer on OpenEarthMap: the transfer does NOT hold there, and why
\newcommand{\cnOemCV}{-0.39}       % five-fold, sd 1.28, 1/5 folds positive
\newcommand{\cnOemCVsd}{1.28}
\newcommand{\cnOemFull}{-0.09}
\newcommand{\cnOemReal}{+1.27}     % land cover DOES improve
\newcommand{\cnOemBg}{-10.94}      % catch-all sacrificed: 17.00 -> 6.06
\newcommand{\cnOemBgAfter}{6.06}
\newcommand{\mrOemBgAfter}{5.83}   % SAM 3 drives it to almost the same value

% ---- third dataset + the supervised-gap framing ----------------------------
\newcommand{\potsBase}{57.83}      % our Potsdam reproduction
\newcommand{\potsPub}{57.8}        % SegEarth-OV3 published -- matches to 0.03
\newcommand{\potsTiles}{2016}      % val tiles, 512^2, stride 512
% ⭐ The oracle is a fully supervised SegFormer-b0 with full training data,
% reported in SegEarth-OV3's own Table 2. Framing the gain as a share of the
% supervised gap is far more legible than a raw delta.
\newcommand{\oracleSup}{50.0}      % LoveDA, fully supervised
\newcommand{\gapClosed}{44\%}      % (48.53-47.37)/(50.0-47.37)

% ---- POTSDAM_RESULTS.md: the third dataset ---------------------------------
\newcommand{\potsShare}{4.29\%}    % `clutter` share of GT -- PRE-REGISTERED
\newcommand{\potsDiscard}{4.68\%}  % measured; committed bracket was 3.78-10.88%
\newcommand{\potsPoint}{4.5\%}     % the committed point estimate
\newcommand{\potsAuc}{0.578}       % best CACHED signal -- prediction was >0.622
% ⚠️ \aucLove is 0.622, LoveDA's best signal INCLUDING texture. For a like-for-like
% comparison with Potsdam (cached signals only) the LoveDA figure is 0.582.
\newcommand{\aucLoveCached}{0.582}
\newcommand{\aucOemCached}{0.794}
\newcommand{\potsReal}{+1.04}      % catch-all-excluded mIoU
\newcommand{\potsFull}{+0.59}      % five-fold full mIoU
\newcommand{\potsFullSd}{0.50}     % spread covers zero -- marginal, say so
\newcommand{\potsDistort}{8.18}    % clutter depresses the headline by this much

% ---------------------------------------------------------------------
%  THE SECOND LEVER: per-class scaling before the argmax.
%  ARGMAX_SCALING_RESULTS.md. All five-fold figures come from one script
%  (argmax_reorder.py) so the substitution table compares like with like.
% ---------------------------------------------------------------------
\newcommand{\scLove}{+1.16}          % LoveDA  C-B, 5-fold
\newcommand{\scLoveSD}{0.19}
\newcommand{\scLoveExcl}{+1.36}      % catch-all-excluded
\newcommand{\scLoveBg}{-0.03}
\newcommand{\scPots}{+4.86}          % Potsdam C-B, 5-fold
\newcommand{\scPotsSD}{0.35}
\newcommand{\scCI}{-0.10}            % ConInfer C-B, 5-fold
\newcommand{\scCISD}{0.14}
\newcommand{\scCIfolds}{1 of 5}
% the substitution table -- all five-fold, same script
\newcommand{\subCItau}{+2.51}
\newcommand{\subLovetau}{+1.16}
\newcommand{\subPotstau}{+0.59}
% end-to-end verification
\newcommand{\eeScaleLove}{+1.34}     % measured increment, LoveDA
\newcommand{\eeScaleLovePred}{+1.31}
\newcommand{\eePotsA}{57.60}
\newcommand{\eePotsB}{58.35}
\newcommand{\eePotsC}{63.27}
\newcommand{\eePotsScale}{+4.92}
\newcommand{\eePotsTotal}{+5.67}
\newcommand{\eePotsMiss}{0.05}       % worst rung miss vs prediction
\newcommand{\eePotsClassMiss}{0.23}  % worst per-class miss
% the tree case
\newcommand{\treeGap}{+54.7}
\newcommand{\treeP}{93.34}
\newcommand{\treeR}{38.63}
\newcommand{\treeTau}{+0.32}         % what per-class tau moved it
\newcommand{\treeScale}{+21.17}      % what the scale moved it, measured
\newcommand{\treeW}{3.73}
\newcommand{\scCalibPots}{100}
\newcommand{\confusionShare}{9.9\%}   % WEEK1 7.6: real-class confusion, LoveDA
\newcommand{\scSpread}{15\%}   % worst real-class fold spread, Potsdam (road)

% ---- the qualitative figure (scripts/fig_qualitative.py, Potsdam val) --------
% Row 4 is a tile where the method LOSES. It is in the figure deliberately: the
% five-fold Potsdam gain \scPots{} is an AVERAGE and four wins would misstate it.
\newcommand{\qualLoss}{7.72}         % 2_13_1024_0_1536_512, full mIoU 41.04->33.32
\newcommand{\qualLossRight}{26.1\%}  % of its changed pixels that became correct

% ---- the open-vocabulary run on Indian UAV imagery (scripts/fig_vocabulary.py)
% ⛔ NO GROUND TRUTH. These are per-class PIXEL SHARES with nodata excluded, not
% accuracy. Nothing derived from them may be phrased as a correctness claim.
\newcommand{\indTiles}{25}           % tiles cut from 5 OpenAerialMap scenes
\newcommand{\indScenes}{5}
\newcommand{\indGSD}{4--5\,cm}
\newcommand{\indFloodA}{78.6}        % `water` under the settlement vocabulary
\newcommand{\indFloodB}{79.1}        % `flood water` under the Indian one
\newcommand{\indBldgA}{18.2}         % `concrete roof`
\newcommand{\indBldgB}{18.3}         % `building`
\newcommand{\indVegA}{100.0}         % solar-farm tile collapses to ONE class
\newcommand{\indSolarB}{4.6}         % ⚠️ PARTIAL -- far below the panels' extent

% the tree case, absolute -- ARGMAX_SCALING_RESULTS.md, measured by eval.py on
% 1816 held-out Potsdam tiles. \treeScale is their difference.
\newcommand{\treeIoUa}{37.92}
\newcommand{\treeIoUb}{59.09}

% ---- combined totals, for the review presentation --------------------------
% Both levers together. LoveDA = \methodGain + \scLove ; Potsdam = eePotsC-eePotsA,
% measured end-to-end by eval.py, not summed from folds.
\newcommand{\loveTotal}{+2.32}       % 1.18 (tau) + 1.16 (scale), 5-fold
\newcommand{\potsTotal}{+5.67}       % 57.60 -> 63.27, eval.py, 1816 held-out tiles
\newcommand{\potsExcl}{+6.00}        % Potsdam catch-all-excluded, 5-fold

% ---- how much of the residual comes back -----------------------------------
% @DISCARD_AFTER_RESULTS.md, 12 Sep. scripts/discard_after.py, LoveDA full 1669
% tiles, 5-fold held out. The accounting identity
%   discard_A - discard_X == recovered - newly discarded
% is verified as exact integers on every run. Rates from a 40k px/tile
% subsample (42.9M real-class px scored); measure_discard_rate.py owns the
% absolutes at the published tau.
\newcommand{\raDiscA}{29.25\%}       % rung A, published tau
\newcommand{\raDiscB}{26.77\%}       % rung B, per-class tau
\newcommand{\raDiscC}{25.01\%}       % rung C, + per-class scale
\newcommand{\raDrop}{4.24}           % points, A -> C
\newcommand{\raRecov}{4.65\%}        % recovered vs A, rung C
\newcommand{\raNew}{0.40\%}          % newly discarded, rung C -- NEVER quote
                                     % \raRecov without this one (WEEK1 8.2)
\newcommand{\raCorrect}{79.2\%}      % of recovered, landing on the right class
\newcommand{\raWPR}{0.26}            % wrong per right, ours
\newcommand{\raTauCorrect}{36.6\%}   % of recovered, tau->0.1 = 1/(1+\tauRatio)
\newcommand{\raFactor}{6.6}          % \tauRatio / \raWPR
% per class, rung C. The four whose tau moved DOWN recover well; the two whose
% tau moved UP recover almost nothing and get it wrong -- by design.
\newcommand{\raWaterA}{31.9\%}
\newcommand{\raWaterC}{18.6\%}
\newcommand{\raWaterOk}{95.0\%}
\newcommand{\raRoadDelta}{+0.1}      % road DISCARDS MORE, deliberately

% ---- separability of the real objective ------------------------------------
% @SEPARABILITY_RESULTS.md. For a FIXED argmax a pixel predicted c either clears
% tau_c or becomes the catch-all -- never another real class -- so TP/FP/FN for
% class c depend on tau_c alone and the `real' objective is a sum of
% single-variable terms. Measured at exactly 0 on every class sweep.
\newcommand{\sepExact}{0.000000}     % max |delta tau| over all sweeps
\newcommand{\sepSweeps}{19}          % class-sweeps, LoveDA + Potsdam + OEM
\newcommand{\sepCostA}{$N\times201$} % cost of exact coordinate ascent
\newcommand{\sepCostB}{$201^{N}$}    % cost of the naive joint search
\newcommand{\sepAllMax}{0.0050}      % Potsdam, --objective all: it does NOT hold

% ---- how much of the oracle the fit captures -------------------------------
% Held-out real-class mIoU. Put this in limitations: it is a measured bound on
% the method, and naming it is stronger than letting a reader compute it.
\newcommand{\capLove}{83\%}
\newcommand{\capPots}{73\%}
\newcommand{\capOEM}{11\%}
\newcommand{\capOemWaterFit}{0.710}   % OEM water, fitted
\newcommand{\capOemWaterOra}{0.240}   % OEM water, oracle
\newcommand{\capOemWaterA}{69.18}     % its IoU at the oracle threshold
\newcommand{\capOemWaterB}{55.02}     % ... and at the fitted one

% ---- the one-standard-error rule: tested, and worse than what we deploy -----
\newcommand{\seLove}{+0.76}          % vs deployed +1.23
\newcommand{\sePots}{+0.55}          % vs deployed +0.90
\newcommand{\seOEM}{+0.54}           % vs deployed +0.28
\newcommand{\seMean}{-0.19}          % mean delta vs the deployed rule

% ---- which classes to calibrate: closed by an UPPER BOUND ------------------
% scripts/freeze_selection.py, LoveDA 1669 tiles, 5-fold held out.
\newcommand{\frzOracle}{+0.08}       % choosing on the EVALUATION fold
\newcommand{\frzCv}{-0.07}           % honest inner-CV selection

% ---- the five levers, one protocol ----------------------------------------
% Levers 1-2 reshape the DECISION and work; 3-5 are reported as bounded
% negatives. Every one was pre-registered.
\newcommand{\lvThreePots}{+0.04}     % head fusion, Potsdam
\newcommand{\lvThreeLove}{+0.22}     % head fusion, LoveDA
\newcommand{\lvFour}{+0.16}          % per-class presence weight, LoveDA
\newcommand{\lvFive}{-0.07}          % per-class bias (vector scaling), LoveDA
\newcommand{\lvFiveSD}{0.18}
\newcommand{\lvRhoCar}{2.33}         % Potsdam car -> instance head
\newcommand{\lvRhoBuild}{0.53}       % Potsdam building -> semantic head
\newcommand{\lvGammaRho}{+0.886}     % Spearman(median S_pres, gamma), p=0.017
\newcommand{\lvGammaP}{0.017}

% ---- sliding-window inference: tested, and it costs -------------------------
% @SLIDING_WINDOW_RESULTS.md. slide_crop=512, stride=341, LoveDA val 1669 tiles.
% SAM 3 resizes input to 1008^2, so 512^2 crops are UPSAMPLED ~2x -- the
% configuration the rest of this literature uses and SegEarth-OV3 does not.
% Precision RISES and recall COLLAPSES, which is a tighter gate rather than
% sharper features: S_pres is computed per view, so a class absent from a crop
% is vetoed there.
\newcommand{\swMIoU}{43.53}          % vs the published single-view 47.38
\newcommand{\swDelta}{-3.85}
\newcommand{\swPrec}{+1.1}           % mPrecision, single-view -> sliding-window
\newcommand{\swRec}{-4.7}            % mRecall
\newcommand{\swWater}{-14.28}        % water IoU, over half the total loss
\newcommand{\swWaterRec}{-16.1}      % water recall
\newcommand{\swBgRec}{+12.9}         % background recall: the catch-all absorbs MORE
\newcommand{\swCrop}{$512$}

% ---- and the gate was MITIGATING the loss, not causing it ------------------
% @SLIDING_WINDOW_RESULTS.md 2a. Loosening the presence gate under sliding
% window costs MORE, and precision collapses while recall barely moves -- so the
% per-crop veto was suppressing false positives, not valid predictions.
\newcommand{\swMax}{39.86}           % gate = max presence over crops
\newcommand{\swGlobal}{39.05}        % gate = one whole-image forward
\newcommand{\swMaxDelta}{-3.67}      % vs the per-crop gate
\newcommand{\swPrecDrop}{68.1 \rightarrow 58.2 \rightarrow 56.7}
\newcommand{\swBuildPrec}{77.0 \rightarrow 49.4}

% ---- dihedral test-time augmentation --------------------------------------
% @TTA_RESULTS.md. Average the score stacks over flipped/rotated views of the
% WHOLE image -- unlike crops, the field of view and therefore the presence
% semantics are identical across views, which is why this works where sliding
% window did not. Aerial imagery has no canonical orientation, so the argument
% is specific to remote sensing.
\newcommand{\ttaHflip}{47.67}        % 2 views, LoveDA val 1669 tiles
\newcommand{\ttaHflipD}{+0.29}
\newcommand{\ttaDfour}{47.84}        % 8 views (all dihedral transforms)
\newcommand{\ttaDfourD}{+0.47}
\newcommand{\ttaPrec}{+0.4}          % mPrecision, hflip -- BOTH rise, uniquely
\newcommand{\ttaRec}{+0.1}           % mRecall
\newcommand{\ttaCost}{6.59}          % s/image at 8 views vs 0.85 baseline
% ⭐ Paired fold-by-fold, same 800 tiles / seed / fold partition. TTA's gain at
% each rung: it survives per-class tau and is ABSORBED by the per-class scale,
% because both change which class wins the argmax.
\newcommand{\ttaAtA}{+0.42}
\newcommand{\ttaAtB}{+0.58}
\newcommand{\ttaAtC}{+0.06}
\newcommand{\ttaScaleWith}{+0.74}    % lever 2's own gain WITH tta
\newcommand{\ttaScaleWithout}{+1.26} % ... and without

% =====================================================================
%  UAVid — the fourth dataset.  @UAVID_RESULTS.md
% =====================================================================
\newcommand{\uavPub}{54.7}           % SegEarth-OV3 Table 2, UAVid^img
\newcommand{\uavBase}{56.86}         % UAVID §1   our reproduction, their vocabulary
\newcommand{\uavGap}{+2.16}          % UAVID §1   unexplained
\newcommand{\uavTiles}{70}           % official val frames
\newcommand{\uavTrainTiles}{200}     % official train frames (real, excl. augmented)
\newcommand{\uavScenes}{20}          % flights in the train split
\newcommand{\uavValScenes}{7}        % flights in the val split
\newcommand{\uavOracle}{59.7}        % their supervised SegFormer-b0 row
\newcommand{\uavCatchShare}{16.16\%} % UAVID §4   catch-all share of GT
\newcommand{\uavDiscard}{6.81\%}     % UAVID §1   real-class pixels to catch-all

% ---- levers on UAVid, group-disjoint folds --------------------------
\newcommand{\uavTauCV}{+1.34}        % UAVID §12  5-fold, 20 flights
\newcommand{\uavTauSD}{0.39}
\newcommand{\uavTauWorst}{+0.75}
\newcommand{\uavScaleCV}{+5.89}      % UAVID §16  5-fold, on top of lever 1
\newcommand{\uavScaleSD}{1.51}
\newcommand{\uavOracleBound}{+1.40}  % UAVID §16  per-class tau oracle, 200 frames

% ---- transfer: fitted on train, applied unchanged to val ------------
\newcommand{\uavTauTrans}{+1.03}     % UAVID §14  lever 1 across the split
\newcommand{\uavTauTransPerm}{0.0\%} % shuffles matching it
\newcommand{\uavScaleTrans}{+5.64}   % UAVID §17  lever 2 across the split
\newcommand{\uavScalePerm}{1.0\%}
\newcommand{\uavScaleShuffle}{-2.82} % UAVID §20  shuffled assignment, corrected prompt
\newcommand{\uavCalibScenes}{3}      % UAVID §14  scenes for +0.81
\newcommand{\uavCalibGain}{+0.81}

% ---- verified end-to-end --------------------------------------------
\newcommand{\uavEEone}{57.92}        % UAVID §15  eval.py, lever 1
\newcommand{\uavEEtwo}{63.55}        % UAVID §18  eval.py, lever 1 + lever 2
\newcommand{\uavEEpred}{63.55}       % cached prediction -- agrees exactly
\newcommand{\uavEEmax}{0.20}         % largest per-class disagreement
\newcommand{\uavaAccA}{79.24}        % UAVID §18  pixel accuracy, baseline
\newcommand{\uavaAccC}{84.81}
\newcommand{\uavPrecD}{+0.60}        % mPrecision, lever 1+2 vs baseline
\newcommand{\uavRecD}{+4.92}         % mRecall

% ---- the corrected configuration -------------------------------------
\newcommand{\uavCorrBase}{60.40}     % UAVID §19  with `low vegetation'
\newcommand{\uavCorrOne}{+1.19}      % lever 1 on the corrected prompt
\newcommand{\uavCorrTwo}{+2.03}      % lever 2 on the corrected prompt
\newcommand{\uavCorrMethod}{+3.22}   % UAVID §20  the honest method contribution
\newcommand{\uavCorrFinal}{63.62}
\newcommand{\uavScaleLost}{64\%}     % share of lever 2's gain the prompt took

% =====================================================================
%  The vocabulary.  @VOCABULARY_RESULTS.md  -- pre-registered, two arms
% =====================================================================
\newcommand{\vocUavGain}{+3.53}      % VOCAB §1   vegetation -> low vegetation
\newcommand{\vocUavTree}{+14.66}     % tree IoU
\newcommand{\vocUavVeg}{+11.14}      % low-vegetation IoU
\newcommand{\vocUavElse}{0.04}       % largest move among the other real classes
\newcommand{\vocUavVegPrecA}{53.62}
\newcommand{\vocUavVegPrecB}{71.51}
\newcommand{\vocUavTreeRecA}{55.90}
\newcommand{\vocUavTreeRecB}{72.76}
\newcommand{\vocPotsBase}{57.83}     % VOCAB §2   recorded Potsdam baseline
\newcommand{\vocPotsNew}{55.12}      % with `low vegetation'
\newcommand{\vocPotsGain}{-2.71}
\newcommand{\vocPotsTreeRec}{38.63}  % IDENTICAL before and after
\newcommand{\vocPotsTreePrecA}{93.34}
\newcommand{\vocPotsTreePrecB}{92.17}
\newcommand{\vocLoveBarren}{+4.94}   % PROMPT_ENSEMBLE, LoveDA `barren'
\newcommand{\vocRouteA}{63.55}       % UAVID §20  bad prompt + both levers
\newcommand{\vocRouteB}{63.62}       % one word + both levers
\newcommand{\vocRouteDelta}{0.07}

% ---- the data traps, reported in setup -------------------------------
\newcommand{\uavLeak}{+0.54}         % UAVID §8   frame-level folds inflate lever 1
\newcommand{\uavLeakA}{+1.05}        % leaking
\newcommand{\uavLeakB}{+0.51}        % group-disjoint
\newcommand{\uavAug}{400}            % augmented copies shipped among 600 `train' frames
\newcommand{\uavFrames}{10}          % consecutive frames per flight

% ---- values that were typed inline until the consistency check caught them ----
% ⚠️ \waterRec is 54.7 and \uavPub is ALSO 54.7 -- water's recall and UAVid's
% published mIoU are unrelated. Never reuse a macro because the value matches.
\newcommand{\tauLove}{0.5}           % published thresholds, per dataset
\newcommand{\tauOEM}{0.1}
\newcommand{\tauPots}{0.1}
\newcommand{\tauUav}{0.3}
\newcommand{\waterPrec}{89.5}        % WEEK1 §5  LoveDA water, the motivating asymmetry
\newcommand{\waterRec}{54.7}
\newcommand{\eeSpreadHi}{0.675}      % WEEK3 §9c fitted span, high end (\eeSpread = low)
\newcommand{\bndWaterPrec}{88.5}     % WEEK3 §9d precision vs oracle tau, non-monotone
\newcommand{\bndRoadPrec}{68.5}
% ⚠️ \oemTauCV is +0.16 and \lvFour is ALSO +0.16 -- OpenEarthMap's threshold gain
% and lever four's null. The second value collision this check caught; matching on
% a value rather than a meaning would have put the wrong citation in the paper.
\newcommand{\oemTauCV}{+0.16}        % WEEK3 §9b OpenEarthMap, 5-fold, full mIoU
\newcommand{\oemTauSD}{0.93}

% =====================================================================
%  DLRSD -- the fifth dataset, and the only one with NO catch-all.
%  @DLRSD_RESULTS.md.  Added 19 Sep 2026; nothing here appears in the
%  paper body yet.
% =====================================================================
\newcommand{\dlrTiles}{2100}         % DLRSD §2  images, 256 x 256
\newcommand{\dlrHeld}{1701}          % DLRSD §1  held-out tiles
\newcommand{\dlrCalib}{399}          % DLRSD §1  calibration draw, stratified
\newcommand{\dlrClasses}{17}
\newcommand{\dlrCatchShare}{0.00\%}  % DLRSD §2  no catch-all class exists
\newcommand{\dlrDiscard}{6.16\%}     % DLRSD §5  real-class pixels to the sink
% ---- the three verified rungs (eval.py, published vocabulary) --------
\newcommand{\dlrEEa}{37.27}          % DLRSD §1  baseline, published tau = 0.1
\newcommand{\dlrEEb}{39.04}          % + per-class tau
\newcommand{\dlrEEc}{44.42}          % + per-class scale
\newcommand{\dlrTotal}{+7.15}
% ---- five-fold, the figures with an error bar -----------------------
\newcommand{\dlrTauCV}{+2.37}        % DLRSD §7  lever 1
\newcommand{\dlrTauSD}{0.63}
\newcommand{\dlrScaleCV}{+5.84}      % DLRSD §8  lever 2 on top of lever 1
\newcommand{\dlrScaleSD}{1.03}
% ---- the oracle: one threshold against seventeen --------------------
\newcommand{\dlrOracleGlobal}{+0.04} % DLRSD §6  best POSSIBLE single tau
\newcommand{\dlrOraclePer}{+2.99}    % best per-class tau
% ---- pixel accuracy, so the mean cannot be dismissed as class averaging
\newcommand{\dlraAccA}{58.94}        % DLRSD §9  eval.py, baseline
\newcommand{\dlraAccC}{65.79}        % both levers
% ---- the vocabulary arm, its own result -----------------------------
\newcommand{\dlrVocBase}{+1.71}      % DLRSD §13 corrected words, before the levers
\newcommand{\dlrVocAfter}{+1.70}     % after them -- the two ADD
\newcommand{\dlrVocEEc}{46.12}       % best verified configuration
% ---- the honest per-class costs -------------------------------------
\newcommand{\dlrWaterLoss}{-2.73}    % DLRSD §9  worst class, under BOTH levers
\newcommand{\dlrTilesUp}{1057}       % DLRSD §10b per-tile outcome
\newcommand{\dlrTilesDown}{769}
\newcommand{\dlrTilesDead}{22}       % emptied entirely -- a stated failure mode
\newcommand{\dlrLeverage}{35.3\%}    % share of the metric owned by classes under 2% of pixels

% =====================================================================
%  The class-scale SEARCH RANGE -- an unreported hyperparameter, measured
%  on four datasets.  @ARGMAX_SCALING_RESULTS.md, 18 Sep 2026.
%  The method's range is 0.40--2.50; these are the sensitivity figures.
% =====================================================================
\newcommand{\rangeLo}{0.40}
\newcommand{\rangeHi}{2.50}
\newcommand{\rangeWideLo}{0.10}
\newcommand{\rangeWideHi}{10}
\newcommand{\rangeDlr}{+1.04}        % wide minus default, DLRSD  (+5.84 -> +6.88)
\newcommand{\rangePots}{+0.66}       % Potsdam (+4.86 -> +5.52)
\newcommand{\rangeLove}{-0.19}       % LoveDA  (+1.16 -> +0.97) and the fit destabilises
\newcommand{\rangeUav}{-0.24}        % UAVid   (+5.89 -> +5.65)
\newcommand{\rangeLoveSDa}{0.19}     % LoveDA fold sd, default range
\newcommand{\rangeLoveSDb}{0.75}     % ... wide range: four times larger
\newcommand{\rangeLoveSpread}{144.7\%} % fitted scales across folds, wide range
\newcommand{\rangeDlrEE}{45.52}      % DLRSD, wide range, verified by eval.py
\newcommand{\rangeDlraAcc}{63.66}    % ... and its pixel accuracy: 65.79 -> 63.66
\newcommand{\rangePotsEE}{63.82}     % Potsdam, wide range, verified


\titleformat{\section}{\large\bfseries\color{acc}}{\thesection.}{0.6em}{}
\newcommand{\why}[1]{\par\smallskip\noindent\textbf{Why this slide exists.} #1}
\newcommand{\say}[1]{\par\smallskip\noindent\textbf{What to say.} #1}
\newcommand{\ask}[2]{\par\smallskip\noindent\textbf{Likely question:} \emph{#1}\\
  \textbf{Answer:} #2}
\newcommand{\care}[1]{\par\smallskip\noindent\bad{\textbf{Be careful.}} #1}

\title{\textbf{Mid-Project Review --- Speaker Notes}\\
\large Per-Class Threshold Calibration for Training-Free Open-Vocabulary\\
Semantic Segmentation of Remote Sensing Imagery}
\author{Priyanshu Kumar \and Raj Gautam}
\date{Companion to \texttt{review.pdf}}

\begin{document}
\maketitle

\noindent\textbf{How to use this.} One entry per slide, in order. Each gives the
purpose of the slide, the line of argument to deliver, and the question most
likely to follow. Total speaking time at this level of detail is roughly
\textbf{18--22 minutes}, which leaves room for discussion in a thirty-minute slot.

\medskip
\noindent\best{The single thread running through the talk:} the pipeline is not
wrong about \emph{what} it sees; it is wrong about \emph{when to commit}. Every
slide either establishes that, quantifies it, or corrects it.

% ---------------------------------------------------------------- 1
\section{Title}
\say{Give the title, the two of you, and the supervisors. One sentence of
orientation: ``This work improves an existing open-vocabulary segmentation
pipeline for aerial imagery without training any network weights.'' Then move
on --- do not explain the method here.}

\section{Presentation Outline}
\say{Read the eight section names in about ten seconds. The purpose is to signal
that the mandated structure is followed, not to summarise.}

% ---------------------------------------------------------------- 2
\section{Introduction}
\why{The committee may not know this sub-field. Three things must land before
anything else: what the task is, what the baseline does, and what we changed.}

\say{``In ordinary semantic segmentation the class list is fixed at training
time. In \emph{open-vocabulary} segmentation the class list is an \emph{input} ---
you supply words at inference, and the model segments them. No weights are
trained; the vocabulary can change from one image to the next.''

\smallskip
``Our baseline is SegEarth-OV3, which applies Meta's SAM~3 to remote sensing. For
each word you give it, SAM~3 returns a score map. The pipeline then takes the
highest-scoring class per pixel, and if that score is below a threshold $\tau$ it
discards the answer and calls the pixel `background'.''

\smallskip
``We reproduced their published result exactly --- \best{47.38 against a
published 47.4} --- so every number after this is measured against a baseline we
control, on our own hardware.''

\smallskip
``Our contribution is to that final decision: one threshold per class, and one
scale per class before the comparison. \best{No weights are trained.}''}

\ask{Why not just train the model?}{Because the setting is training-free by
definition: the class list is not known in advance, so there is nothing fixed to
train for. Fine-tuning also needs labelled data in the thousands of tiles, which
is exactly what a new region does not have.}

\care{Do not say ``unsupervised''. The vocabulary is given, so the correct terms
are \emph{training-free} and \emph{annotation-free}.}

% ---------------------------------------------------------------- 3
\section{Motivation --- the baseline discards land cover}
\why{This is the slide that justifies the entire project. If the committee
accepts that a large, recoverable residual exists and that no simple fix reaches
it, everything downstream follows.}

\say{``We instrumented the baseline to ask a question its own paper does not:
how much real land cover does it throw away? On LoveDA, at their own published
threshold, \best{29.68\% of all real-class pixels} --- 323 million --- are
assigned to the catch-all class.''

\smallskip
``More precisely: the pipeline is not mostly \emph{wrong}, it is mostly
\emph{silent}. Discard exceeds genuine class confusion by three to one.''

\smallskip
``The clearest case is \texttt{water}. Precision \best{89.5\%}, recall
\best{54.7\%}. When it says water it is right nine times out of ten --- and it
finds barely half the water in the image. Those missing pixels are not
misclassified. They fall below the threshold and are discarded.''

\smallskip
``The obvious response is to lower the threshold. We tried it: \bad{$-5.54$
mIoU}, because you gain 1.73 wrong pixels for every right one. We also tried
removing the presence gate: \bad{$-11.97$}. Both global adjustments cost more
than they return.''

\smallskip
``So the correction cannot be global. \best{It has to act per class} --- which is
the rest of the talk.''}

\ask{Isn't a low recall just the model being bad at water?}{No, and the precision
number is why. A model that is genuinely bad at a class is wrong when it fires.
This one is right 89.5\% of the time it fires. The information is present; the
decision rule refuses to use it.}

% ---------------------------------------------------------------- 4
\section{Motivation --- the discard, visualised}
\why{The previous slide is numbers. This is the same fact as a picture, and it is
what people remember.}

\say{``Four panels, one LoveDA tile. Input, ground truth, the baseline's
prediction, and in the last panel the pixels it sent to background.''

\smallskip
``Look at the ground truth: the top of the tile is water. Look at the
prediction: that region is gone. And in the last panel you can see it was
\emph{discarded}, not confused with another class --- \best{19\% of this single
tile}.''

\smallskip
``The threshold that caused this was tuned for the dataset as a whole. It is a
reasonable value on average and the wrong value for water.''}

\ask{Is this tile representative?}{It is one tile chosen to show the mechanism
clearly. The dataset-wide figure is the 29.68\% on the previous slide, measured
over all 1669 validation images.}

% ---------------------------------------------------------------- 5
\section{Application domains}
\why{Establishes that the training-free setting is not an academic convenience
but the operational reality in several domains.}

\say{``Where does an open vocabulary actually matter? Three cases.''

\smallskip
``\textbf{Disaster response.} After a flood you have UAV imagery within hours and
no labelled data for that district. You cannot train anything in time.''

\smallskip
``\textbf{Infrastructure.} A class like \texttt{solar panel} appears in no
land-cover benchmark. Here you can simply type it.''

\smallskip
``\textbf{Settlement mapping.} The vocabulary changes between surveys ---
\texttt{tin roof}, \texttt{unpaved road}, \texttt{water tank}.''

\smallskip
``The relevant number is the deployment cost: our correction needs
\best{about \calibTiles{} labelled tiles and no training run}. Fine-tuning
instead would need thousands of labelled tiles and would forfeit the
training-free property that made the pipeline attractive.''}

% ---------------------------------------------------------------- 6
\section{Problem statement --- formulation}
\why{The committee asked for a mathematical formulation. This slide is that, and
it should take under two minutes.}

\say{``$s_c(p)$ is the score class $c$ gets at pixel $p$. The baseline takes the
\texttt{argmax} --- the highest-scoring class --- and keeps it only if that score
clears one global $\tau$; otherwise the pixel becomes the catch-all class $b$.''

\smallskip
``\texttt{argmax} simply means `which class scored highest'. \texttt{max} gives
the value, \texttt{argmax} gives which one it was.''

\smallskip
``We change two things. First, $\tau$ becomes a \emph{vector}: each class gets
its own threshold. Second, we multiply each class's score by a fitted weight
$w_c$ \emph{before} the comparison, so the weight can change \emph{which} class
wins.''

\smallskip
``Both are fitted by coordinate ascent on about \calibTiles{} labelled tiles, and
--- this is the point --- \best{no network weights are modified}. The threshold
is applied to the raw score, so $w$ only affects the \texttt{argmax}.''}

\ask{Isn't fitting on labelled data just supervision through the back door?}{The
baseline already tunes its single $\tau$ on labelled data --- 0.5 for LoveDA, 0.1
for OpenEarthMap. We spend the same kind of budget on $N$ parameters instead of
one. It is the same protocol, not a new one.}

% ---------------------------------------------------------------- 7
\section{Completeness of the per-class threshold family}
\why{\best{This is the intellectual core of the work.} If a committee member asks
``what is actually novel here?'', this slide is the answer. Do not rush it.}

\say{``A natural objection is: why thresholds? Why not recalibrate the scores in
some cleverer way?''

\smallskip
``The proposition answers that. Take \emph{any} strictly increasing per-class
transformation of the scores, and compare against a global threshold. Because the
transformation is increasing, you can invert it --- and the comparison becomes
exactly a per-class threshold.''

\smallskip
``So \best{per-class thresholds are not one option among many. They are the
complete family} of everything reachable by rescaling scores \emph{after} the
\texttt{argmax}. Nothing in that family can beat them.''

\smallskip
``Scaling \emph{before} the \texttt{argmax} is different, because it changes
which class wins. That is a strictly larger family --- and it is not decoration.
We measured that \best{6\% of discarded pixels are cases where the catch-all
class won the \texttt{argmax} outright}. A threshold cannot change an
\texttt{argmax}, so no threshold anywhere can fix those. That is precisely what
the second parameter set is for.''}

\ask{So the second lever subsumes the first --- why keep both?}{They are fitted
together and reported as two stages so the contribution of each is visible. The
first stage is also the one with the completeness guarantee, which is worth
stating separately.}

% ---------------------------------------------------------------- 8
\section{Literature survey}
\why{Committee requirement: a minimum of ten works. Fourteen are listed.}

\say{Do not read the table. Walk the three groups in about a minute:

\begin{itemize}[nosep,leftmargin=1.4em]
  \item \textbf{Rows 1--7} --- how dense prediction was extracted from
        vision--language models: CLIP, then MaskCLIP removing pooling, then a
        line of attention corrections.
  \item \textbf{Rows 8--11} --- the remote-sensing line, ending at our baseline
        SegEarth-OV3, and ConInfer as the nearest contemporary work.
  \item \textbf{Rows 12--14} --- the tools we borrow: calibration, classical
        thresholding, classical spatial refinement.
\end{itemize}}

\care{You have read SegEarth-OV3, ConInfer and SAM~3 in full. If asked about a
row you have only skimmed, say so plainly and describe what it contributes ---
that is a normal and acceptable answer.}

% ---------------------------------------------------------------- 9
\section{Research gaps}
\why{Connects the literature to the contribution. Each gap should map to
something the work then does.}

\say{``\textbf{Gap one.} Every method in this family ends the same way: an
\texttt{argmax}, then one global threshold. All the research effort goes into
producing better \emph{scores}. Nobody has asked whether the \emph{rule} applied
to those scores is the right shape.''

\smallskip
``\textbf{Gap two.} Papers report mIoU. None reports how much real land cover
ends up in the catch-all class, or whether it can be recovered. We measured it.''

\smallskip
``\textbf{Gap three.} mIoU averages classes \emph{equally}. So a `background'
class --- which in LoveDA means `everything else' --- owns one seventh of the
score regardless of whether it means anything. A method can appear to improve or
degrade while land-cover quality is unchanged. We observed this in both
directions and therefore report two metrics throughout.''}

% ---------------------------------------------------------------- 10
\section{Objectives}
\why{Separates what is finished from what is not. \best{Do not present the
project as complete.}}

\say{``Three objectives are met. Reproduce the baseline exactly, so our numbers
mean something. Quantify what it discards and test whether any global
adjustment recovers it --- none does. And design a correction that trains no
weights.''

\smallskip
``Three remain open. The fitted parameters are currently \emph{constant across
tiles}, so they are right on average and can be wrong on an individual scene; a
scene-adaptive version is the natural next step. SAM~3's two heads are combined
by a fixed maximum, which is another unexamined decision of the same kind. And
the evaluation needs widening to establish when the correction transfers.''}

\ask{What is the biggest remaining risk?}{Transfer. The fitted parameters are
specific to a data distribution, so a deployment must calibrate on the
distribution it will run on. Characterising that boundary properly is the main
open item.}

% ---------------------------------------------------------------- 11
\section{Experimental setup}
\why{Establishes that the evaluation is complete and honest before any result is
shown.}

\say{``Three datasets, and \best{every image in each is used}. LoveDA
validation is 1669 tiles at 0.3\,m; Potsdam is 2016 tiles at 5\,cm;
OpenEarthMap is \oemTiles{} tiles.''

\smallskip
``Two points about the protocol matter. \best{Calibration and evaluation tiles
are always disjoint} --- we never fit and test on the same image. And the
baseline is \emph{recomputed} on the same held-out tiles, so the two rules are
compared on identical pixels rather than against a published number.''

\smallskip
``We also keep a validation gate: any instrumented run must still reproduce
47.37 mIoU and 29.68\% discard. If either moves, the instrumentation changed the
model and the run is discarded.''}

\ask{Why only \oemTiles{} tiles of OpenEarthMap?}{That is the public
redistribution we could obtain; the official split is 500. We say so wherever
that dataset appears.}

% ---------------------------------------------------------------- 12
\section{Results}
\why{The main quantitative claim.}

\say{``Read the table by column. On LoveDA the two stages give
\best{\methodGain{}} and \best{\scLove{}}, total \best{\loveTotal{}}. On Potsdam,
\best{\potsTotal{}}. On OpenEarthMap the headline is $+0.30$ --- but under the
catch-all-excluded metric it is \best{\mrOemReal{}}, larger than LoveDA's. That
is gap three from earlier, appearing in our own results.''

\smallskip
``\best{Every fold is positive} on both datasets where we ran the full protocol.''

\smallskip
``The row that matters most is the verification. These gains were first computed
from a cached statistic. We then ran the actual pipeline end to end on Potsdam:
\best{\eePotsA{} $\rightarrow$ \eePotsB{} $\rightarrow$ \eePotsC{}} on 1816
held-out tiles, \best{every stage within 0.05 mIoU} of what we predicted.''

\smallskip
``The largest single change is Potsdam \texttt{tree}: \best{\treeIoUa{} to
\treeIoUb{}}. Its precision is \treeP{}\% and its recall \treeR{}\% --- the same
asymmetry as \texttt{water} on LoveDA, and exactly what the method is built to
correct.''}

\ask{Is a one-to-two point gain significant?}{Five-fold cross-validation with
every fold positive, and the Potsdam gain is \potsTotal{}. For comparison, the
gap between the baseline and the previous state of the art in this family is of
the same order --- and this costs no training.}

% ---------------------------------------------------------------- 13
\section{Results --- qualitative}
\why{Shows the correction acting, and shows a failure. \best{Point at the failure
before anyone else does.}}

\say{``Four Potsdam tiles. Input, ground truth, baseline, ours. \best{The network,
its weights and the forward pass are identical between the last two columns} ---
only the decision rule differs.''

\smallskip
``Row one: a normal urban scene. \texttt{tree} improves, and the classes that
were already well calibrated --- \texttt{building}, \texttt{car} --- are left
within a point of where they were. The correction does not churn.''

\smallskip
``Row two: the baseline misses a canopy entirely; the correction recovers it.''

\smallskip
``Row three: fourteen pixels change out of a quarter of a million. Where nothing
is miscalibrated, nothing happens.''

\smallskip
``\bad{Row four is a tile where we lose \qualLoss{} mIoU.} It is in the figure
deliberately. Our parameters are constant across tiles, so they are correct on
average --- \potsTotal{} over the dataset --- and can be wrong on one scene. That
is the limitation, and it is the open objective on the previous slide.''}

% ---------------------------------------------------------------- 14
\section{Application study}
\why{Demonstrates the method outside the benchmarks, on imagery from India.}

\say{``\indTiles{} tiles from \indScenes{} openly licensed UAV scenes over India
at \indGSD{} --- a region in none of the three benchmarks. We ran it unchanged
under two different class lists, typed at inference.''

\smallskip
``The finding is the \emph{stability}. A flooded building is partitioned almost
identically by both vocabularies: \indFloodA{}\% under \texttt{water} against
\indFloodB{}\% under \texttt{flood water}; \indBldgA{}\% under
\texttt{concrete roof} against \indBldgB{}\% under \texttt{building}. \best{The
same regions are recovered; only the names follow the words supplied.}''

\smallskip
``The second case is a failure and worth stating. A solar farm: with no word for
panels, \indVegA{}\% of the tile goes to one class. Supply \texttt{solar panel}
and it recovers \indSolarB{}\% --- \bad{below the panels' visible extent}.''

\smallskip
``\bad{There is no ground truth for this imagery}, so these are pixel shares, not
accuracy. We make no accuracy claim on these scenes.''}

\ask{Why not calibrate on the Indian imagery too?}{It is unlabelled, and our
transfer results show that parameters fitted elsewhere would not help. The
\calibTiles{} labelled tiles are what such a deployment would have to budget for
--- which is the honest cost of the method.}

% ---------------------------------------------------------------- 15
\section{Application study --- one scene, two vocabularies}
\say{``Same weights, same forward pass. Only the class list differs between the
second and third columns. Lower row, the flooded building --- recovered
consistently under both. Upper row, the photovoltaic array, for which the first
vocabulary has no corresponding class at all.''}

% ---------------------------------------------------------------- 16
\section{Project timeline}
\say{``Work to date: the baseline and the measurement in August; the two
parameter sets and their verification through September; the application study
this month. Also in August we built a semantic co-occurrence prior, evaluated it
and \emph{rejected} it --- it did not earn its place, and that is recorded.''

\smallskip
``Remaining: the scene-adaptive formulation and the head-fusion question through
December, then writing, with submission targeted at a journal in the first
quarter of 2027.''}

\ask{Why was the co-occurrence prior dropped?}{We measured it against a plain
neighbour vote. It added two tenths of a point, and a \emph{perfect} matrix added
three tenths. The idea was sound and the measurement did not support it, so it is
reported as an ablation rather than a contribution.}

% ---------------------------------------------------------------- 17
\section{References and closing}
\say{Do not read them. ``Fourteen key references; the full bibliography is in the
manuscript.'' Then thank the committee and invite questions.}

\bigskip\hrule\bigskip
\section*{Three things to have ready}
\begin{enumerate}[leftmargin=1.6em]
  \item \best{``What is novel?''} --- the completeness proposition. Per-class
        thresholds are the complete family given a fixed \texttt{argmax}; scaling
        before the \texttt{argmax} is provably larger, and reaches the 6\% of
        errors a threshold cannot.
  \item \best{``Is this just a SAM~3 quirk?''} --- no. The same correction was
        applied to a CLIP-based pipeline, where it is worth \cnCV{}$\pm$\cnCVsd{}~mIoU,
        larger than on SAM~3. One architecture is a result; two is a property.
  \item \best{``What is left?''} --- scene-adaptive parameters, the head-fusion
        question, and characterising when the correction transfers. Say this
        plainly; the work is not finished.
\end{enumerate}

\end{document}
